\chapter{The neural baselines}
\label{ch:neural-baselines}

The paper compares its structured estimator with six neural sequence denoisers. This
chapter explains each of them from scratch, for a reader who has never built a neural
network: what a layer is, what the six architectures compute, how many numbers each one
holds, how the noise level enters, what the network is asked to output and why that
output is a score, the single training recipe every architecture shares, the two losses
(denoising and implicit score matching), the three ways of computing the divergence
that the implicit loss needs, and what experiments E03, E08 and E14 do with all of this.
Everything here is implemented in one file, \code{src/markov\_score/nn\_torch.py}, and
the section headings follow it.

\section{From a map to a score}
\label{sec:net-io}

\subsection{What a denoiser network is}

A \term{neural network} here is a function $f_\theta:\Reals^{\nsites}\times(0,\infty)\to\Reals^{\nsites}$
that takes a noised sequence $\vx$ and a level $t$ and returns one number per site,
built from a fixed sequence of simple operations whose numerical constants, the
\term{parameters} $\theta$, are chosen by minimising a loss on data. The paper's networks
are \term{denoisers}: their output is used to form an estimate of the posterior mean
$\pmean{\va}=\E[\va\mid\vx_t=\vx]$, from which Tweedie's identity
(\Cref{ch:ou-channel-tweedie}) gives an estimate of the score,
\begin{equation}
  \hat\score(\vx,t)=-\frac{\vx-e^{-t}\hat m(\vx,t)}{\Dt}.
  \label{eq:score-from-mean}
\end{equation}
The paper's estimand compares $\hat\score$ with the tree-exact score, so the network is
judged as a score estimator whatever it was trained to output.

\subsection{The two output parameterisations}

Write $\alpha=e^{-t}$ and $\vx=\alpha\va+\sqrt{\Dt}\vz$. Two ways of turning the
network's output into $\hat m$ are kept, because the diffusion literature uses both and
the screen lets each family pick the one it trains better with:
\begin{itemize}
\item the \term{x0 parameterisation}: the output is the mean itself, $\hat m=f_\theta(\vx,t)$;
\item the \term{eps parameterisation}: the output is an estimate $\hat\varepsilon$ of the
  standardised noise $\vz$, and the mean is recovered by inverting the channel,
  $\hat m=(\vx-\sqrt{\Dt}\,\hat\varepsilon)/\alpha$.
\end{itemize}
(\code{Denoiser.posterior\_mean}.) At the optimum the two are equivalent, because
$\E[\vz\mid\vx]=(\vx-\alpha\E[\va\mid\vx])/\sqrt{\Dt}$ is a bijective affine image of
$\E[\va\mid\vx]$; they differ in how the training loss weights the levels and in the
scale of the output the network must produce (order one for eps at every $t$; order one
for x0 too, but with a factor $1/\alpha$ amplifying its errors in the mean at high $t$).
Substituting the eps form into \cref{eq:score-from-mean} gives the identity
$\hat\score=-\hat\varepsilon/\sqrt{\Dt}$: an eps network is a score network up to the
scalar $-1/\sqrt{\Dt}$.

\subsection{The time features}

The level enters every architecture through a small vector, the \term{Fourier features}
of $\log t$ (\code{TimeFeatures}):
\begin{equation}
  \tau(t)=\bigl(\log t,\ \sin(2^k\log t),\ \cos(2^k\log t)\bigr)_{k=0,1,2,3}\in\Reals^{9}.
  \label{eq:time-features}
\end{equation}
Two choices are deliberate. The argument is $\log t$ rather than $t$ because the
twelve levels of the schedule are log-spaced on $[0.05,3]$ (\Cref{ch:protocol}), so in
$\log t$ they are equally spaced and the network sees a uniform grid. The sines and
cosines at four octaves give the network smooth, non-monotone basis functions of the
level, so that a level-dependent gain (the correct denoiser shrinks less at low noise
and more at high) can be represented by a linear combination in the first layer. The
MLP concatenates $\tau(t)$ once to its input; every other architecture projects it with
a linear map to its channel width and adds the result at every site, so that the level
modulates every position identically, which is the natural symmetry of a coordinatewise
channel.

\section{The building blocks}
\label{sec:blocks}

\paragraph{Affine layers and activations.}
An \term{affine layer} (\code{nn.Linear}) maps $\vect{h}\in\Reals^{n}$ to
$\mat{W}\vect{h}+\vect{b}$ with a weight matrix $\mat{W}\in\Reals^{m\times n}$ and a bias
$\vect{b}\in\Reals^m$: $mn+m$ parameters. An \term{activation} is a fixed nonlinear
scalar function applied entrywise; the code uses $\tanh$ and the \term{GELU},
$x\,\Phi(x)$ with $\Phi$ the standard normal distribution function, a smooth relative of
the rectifier. Without activations a stack of affine layers is one affine layer, so the
activations are what make a network more than a linear map.

\paragraph{One-dimensional convolution.}
A \term{1-D convolution} with $C$ input channels, $O$ output channels, kernel width $k$
and dilation $d$ maps a signal $\vect{h}\in\Reals^{C\times\nsites}$ to
\begin{equation}
  (\mat{W}\ast\vect{h})_{o,i}=b_o+\sum_{c=1}^{C}\sum_{j=0}^{k-1}W_{o,c,j}\,h_{c,\,i+d(j-r)},
  \qquad r=\tfrac{k-1}2,
  \label{eq:conv}
\end{equation}
with $h_{c,i}=0$ for $i$ outside $\{1,\dots,\nsites\}$ (\term{zero padding}). It has
$OCk+O$ parameters, the same weights at every site $i$ (\term{weight sharing}), and its
output at $i$ depends on the inputs at sites $i-dr,\dots,i+dr$ only: its
\term{receptive radius} is $dr$. Weight sharing is the network's version of the chain's
homogeneity (the same kernel at every transition), and the receptive radius is the
network's version of locality. The code implements \cref{eq:conv} as a gather followed by
a matrix product (\code{FastConv1d}), which is algebraically identical to the framework's
convolution and was ten to a hundred times faster on the laptop's CPU; the test suite
checks the two agree to round-off.

\paragraph{Residual connections and layer normalisation.}
A \term{residual connection} replaces a block $g$ by $\vect{h}\mapsto\vect{h}+g(\vect{h})$,
so that the block learns a correction to the identity and deep stacks train stably.
\term{Layer normalisation} rescales a vector to zero mean and unit variance across its
channels and then applies a learned per-channel gain and shift ($2d$ parameters for
width $d$); the transformer applies it before each block (\term{pre-LN}).

\section{The six architectures}
\label{sec:archs}

\subsection{The fully connected network (MLP)}

The \term{multilayer perceptron} sees the whole sequence at once, with no structure
imposed: the input is the concatenation $(\vx,\tau(t))\in\Reals^{\nsites+9}$, followed by
two hidden affine layers with $\tanh$ and an output affine layer to $\Reals^{\nsites}$:
\begin{equation}
  f(\vx,t)=\mat{W}_3\tanh\bigl(\mat{W}_2\tanh(\mat{W}_1(\vx,\tau(t))+\vect{b}_1)+\vect{b}_2\bigr)+\vect{b}_3 .
\end{equation}
Every output depends on every input through dense matrices, there is no weight sharing
across sites, and the network is tied to the length it was trained at (\code{length\_agnostic
= False}), so it is excluded from the length study. It is the baseline against which
weight sharing is measured.

\subsection{The weight-shared window head}

The \term{window head} is the smallest local architecture: one convolution of width
$2r+1$ from one channel to $w$ channels (zero-padded, so the output has the input's
length), the level added through a linear projection of $\tau(t)$ to the $w$ channels,
then a \term{pointwise MLP}, that is, two convolutions of width one (affine maps applied
at each site separately) with $\tanh$ between and after the first:
\begin{equation}
  \vect{h}_i=\tanh\bigl(\mat{W}_{\mathrm{conv}}\ast\vx\big|_i+\mat{W}_\tau\tau(t)+\vect{b}\bigr),\qquad
  f(\vx,t)_i=\vect{w}_2^\top\tanh(\mat{W}_1\vect{h}_i+\vect{b}_1)+b_2 .
\end{equation}
Its output at site $i$ depends on $x_{i-r},\dots,x_{i+r}$ only: the receptive radius is
exactly $r$, and the module declares it (\code{receptive\_radius = radius}), which is what
the coloured-probe divergence of \Cref{sec:divergence} relies on. It is length-agnostic,
since nothing in it refers to $\nsites$. The screen chose radius $4$ and width $64$ for
E03; E08 uses radius $2$ and width $64$.

\subsection{The dilated temporal convolutional network (TCN)}

A \term{temporal convolutional network} stacks convolutions of width three with
dilations $d\in\{1,2,4,8\}$, each in a residual block:
\begin{equation}
  \vect{h}^{(0)}=\mat{W}_0\ast\vx+\mat{W}_\tau\tau(t),\qquad
  \vect{h}^{(\ell)}=\vect{h}^{(\ell-1)}+\mat{P}_\ell\bigl(\tanh(\mat{W}_\ell\ast_{d_\ell}\vect{h}^{(\ell-1)})\bigr),\qquad
  f=\vect{w}^\top\vect{h}^{(4)}+b,
\end{equation}
with $\mat{P}_\ell$ a width-one convolution and $\ast_{d}$ the dilated convolution
\cref{eq:conv}. The convolutions are non-causal (they look both ways), as they should
for a smoother. Each block of dilation $d$ widens the receptive radius by $d$, so the
stack sees $1+2+4+8=15$ sites on either side: a long range at a cost that grows linearly
in the depth rather than in the radius, which is the point of dilation. It is
length-agnostic. The screen chose $32$ channels.

\subsection{The bidirectional gated recurrent unit (bi-GRU)}

A \term{recurrent network} processes the sequence site by site, carrying a hidden state
$\vect{h}_i\in\Reals^H$ from $i-1$ to $i$. The \term{gated recurrent unit} updates it by
\begin{equation}
\begin{aligned}
  \vect{r}_i&=\sigma(\mat{W}_r\vect{u}_i+\mat{U}_r\vect{h}_{i-1}+\vect{b}_r),\qquad
  \vect{z}_i=\sigma(\mat{W}_z\vect{u}_i+\mat{U}_z\vect{h}_{i-1}+\vect{b}_z),\\
  \tilde{\vect{h}}_i&=\tanh\bigl(\mat{W}_n\vect{u}_i+\vect{r}_i\odot(\mat{U}_n\vect{h}_{i-1}+\vect{b}_n)+\vect{b}_n'\bigr),\qquad
  \vect{h}_i=(1-\vect{z}_i)\odot\tilde{\vect{h}}_i+\vect{z}_i\odot\vect{h}_{i-1},
\end{aligned}
\end{equation}
where $\sigma$ is the logistic function, $\odot$ the entrywise product, and
$\vect{u}_i=\tanh(\mat{W}_{\mathrm{in}}(x_i,\tau(t))+\vect{b}_{\mathrm{in}})$ the per-site
input. The \term{reset gate} $\vect{r}_i$ and the \term{update gate} $\vect{z}_i$ let the
state be overwritten or kept, which is what allows information to travel far along the
chain without vanishing. \term{Bidirectional} means two such recurrences, one running
left to right and one right to left, whose states are concatenated; a per-site head then
maps $(\vect{h}_i^{\rightarrow},\vect{h}_i^{\leftarrow},x_i)$ through one hidden layer to the
output. The analogy with the forward and backward messages of \Cref{ch:factor-graphs-bp}
is exact in shape and inexact in content: the GRU's state is a learned vector, not a
function on the grid. The recurrence has $6H^2+6H$ parameters per direction. Global
propagation, shared weights, length-agnostic; the screen chose $H=32$.

\subsection{The transformer encoder with relative-position bias}

The \term{transformer} replaces recurrence by \term{self-attention}: every site forms a
weighted average of every other site's representation, with weights it computes
itself. With $\mat{H}\in\Reals^{\nsites\times d}$ the current representation (row $i$ for
site $i$), one head computes queries, keys and values $\mat{Q}=\mat{H}\mat{W}_q$,
$\mat{K}=\mat{H}\mat{W}_k$, $\mat{V}=\mat{H}\mat{W}_v$ of width $d_h=d/\text{heads}$ and
\begin{equation}
  \mathrm{att}_{ij}=\frac{\exp\bigl(\vect{q}_i\cdot\vect{k}_j/\sqrt{d_h}+\beta_{c(i-j)}\bigr)}{\sum_{j'}\exp\bigl(\vect{q}_i\cdot\vect{k}_{j'}/\sqrt{d_h}+\beta_{c(i-j')}\bigr)},
  \qquad
  \mathrm{out}_i=\sum_j\mathrm{att}_{ij}\,\vect{v}_j ,
  \label{eq:attention}
\end{equation}
where $c(k)=\max(-D,\min(D,k))$ clips the offset at $D=32$ and $\beta_{-D},\dots,\beta_D$
are $2D+1$ learned numbers per head, the \term{relative-position bias}
(\code{RelPosAttention}). The heads are concatenated and projected; the block is
$\mat{H}\leftarrow\mat{H}+\mathrm{Attn}(\mathrm{LN}(\mat{H}))$ followed by
$\mat{H}\leftarrow\mat{H}+\mathrm{FF}(\mathrm{LN}(\mat{H}))$ with a two-layer feed-forward
map of inner width $2d$ and GELU, repeated for two layers; a final layer normalisation
and a linear map give one output per site. Why this design is \emph{length-agnostic}:
the network never sees an absolute position. It sees $x_i$, the level, and offsets
$i-j$ through the bias, so the same weights apply to a sequence of any length. Why it
nevertheless fails to extrapolate in the paper's length study is also visible in
\cref{eq:attention}: with training length $\nsites=32$ the offsets range over
$-31,\dots,31$, so the clipped entries $\beta_{\pm32}$ never receive a gradient and stay
at their zero initialisation, while at test length $256$ most pairs have $|i-j|\ge32$
and use exactly those entries; and the softmax normalises over eight times more keys,
which changes every attention weight. The paper attributes the degradation to the
positional encoding rather than to attention as such, with references to the length
extrapolation literature. The screen chose $d=32$, four heads, two layers.

\subsection{The two-level one-dimensional U-Net}

A \term{U-Net} processes the sequence at two resolutions. An encoder of two width-three
convolutions with GELU acts at full resolution ($c$ channels); a strided convolution
(width four, stride two) halves the length and doubles the channels; a second encoder
acts at half resolution; a transposed convolution (the adjoint of a strided
convolution) restores the length; the result is concatenated with the full-resolution
encoder output (the \term{skip connection}) and a decoder of two more convolutions and a
width-one head produce the output. The level is added after the first encoder. The
coarse branch sees a receptive field twice as wide per layer, so the network combines
local detail with a wider, cheaper view; the skip connection keeps the fine detail the
downsampling would lose. Multiscale, shared weights, length-agnostic (a one-site
padding fixes odd lengths); the screen chose $c=16$.

\subsection{Parameter counts}

The number of parameters of a network is fixed by its architecture and can be counted
by hand; the code's \code{n\_params} sums them. For the six screened configurations the
counts are $26{,}016$ (MLP), $5{,}505$ (window head), $17{,}057$ (TCN), $15{,}169$ (bi-GRU),
$18{,}057$ (transformer) and $13{,}713$ (U-Net); these are the values in the paper's screen
table, obtained there from the code and here from the architecture, which is a useful
check that the two agree.

\begin{example}[Two parameter counts, by hand]
\label{ex:params}
\emph{The window head} at radius $4$, width $64$: the first convolution has one input
channel, $64$ output channels and width $2\cdot4+1=9$, so $64\cdot1\cdot9+64=640$
parameters; the level projection is an affine map from $9$ to $64$, $9\cdot64+64=640$;
the pointwise MLP is a width-one convolution from $64$ to $64$, $64\cdot64+64=4160$, and
one from $64$ to $1$, $64+1=65$. Total $640+640+4160+65=5505$. \emph{The MLP} with two
hidden layers of $128$ at $\nsites=32$: the input has $32+9=41$ entries, so the layers hold
$41\cdot128+128=5376$, $128\cdot128+128=16512$ and $128\cdot32+32=4128$ parameters,
total $26016$. The MLP has five times the parameters of the window head and, at every
sample size in the paper's baseline table, several times its error: what it lacks is
not capacity but the weight sharing and locality that match the data.
\end{example}

\section{The one training recipe}
\label{sec:recipe}

Every architecture, at every interface, trains under the same \code{TrainConfig}:
\begin{itemize}
\item \term{batch} $64$: each gradient step uses $64$ sequences drawn with replacement
  from the training set;
\item \term{per-example level}: each sequence in the batch receives its own level $t$,
  drawn uniformly from the twelve frozen levels, so a batch mixes levels and the
  network learns all of them together;
\item fresh \term{corruption} at every step under interface~A: $\vz\sim\normal0{\Id}$ is
  redrawn for every sequence in every batch (\code{\_draw\_batch});
\item \term{Adam}, a first-order optimiser that scales each parameter's step by a
  running estimate of its gradient's magnitude, with learning rate $10^{-3}$ or
  $3\cdot10^{-3}$ (the screen chooses);
\item \term{gradient clipping} at norm $1$: if the gradient of the loss over all
  parameters has Euclidean norm above $1$ it is rescaled to norm $1$ before the step,
  which prevents a rare large batch from throwing the parameters far;
\item a \term{checkpoint ladder} $\{500,1000,2000,4000,8000,16000,32000\}$: the parameters
  are copied at each rung and the rung is chosen afterwards (\Cref{ch:model-selection});
\item both output parameterisations screened, one run per configuration in the screen
  (the pre-data amendment of E03 cut three restarts to one for budget), and a seed for
  the initialisation and the batch stream derived from the cell's tags by
  \code{rng\_for}, so a run is reproducible from its row;
\item one thread per process; the confirmatory runs used the CPU code path
  operation for operation.
\end{itemize}
The recipe is what makes a difference between two rows of the baseline table a
difference between architectures. The earlier screen of this model, which the audit
found had given the long-range architectures one level per step and one restart against
the window head's twelve and eight, is the counterexample the recipe was written against.

\section{The two losses}
\label{sec:losses}

\subsection{Denoising score matching (DSM)}

Under interface~A the network regresses on a clean target:
\begin{equation}
  \mathcal L_{\mathrm{DSM}}(\theta)=\E\bigl\lVert f_\theta(\vx,t)-\vect{y}\bigr\rVert^2,
  \qquad
  \vect{y}=\begin{cases}\vz&\text{(eps)},\\ \va&\text{(x0)},\end{cases}
  \label{eq:dsm}
\end{equation}
averaged over the batch, the level and the corruption (\code{dsm\_loss}). Why this
learns the score is a two-line argument. For a fixed $(\vx,t)$ the squared loss is
minimised by the conditional expectation of the target,
$f^\ast(\vx,t)=\E[\vect{y}\mid\vx,t]$, because $\E\lVert f-\vect{y}\rVert^2=\E\lVert f-\E[\vect{y}\mid\vx]\rVert^2+\E\lVert\vect{y}-\E[\vect{y}\mid\vx]\rVert^2$
and only the first term depends on $f$. For the x0 target this is $\E[\va\mid\vx]$, the
posterior mean, and \cref{eq:score-from-mean} turns it into the score. For the eps
target, $\E[\vz\mid\vx]=(\vx-\alpha\E[\va\mid\vx])/\sqrt{\Dt}=-\sqrt{\Dt}\,\score(\vx,t)$:
the optimal eps network is minus $\sqrt{\Dt}$ times the score. This is Vincent's
connection between denoising and score matching, and \Cref{ch:score-matching} gives it
with more care. The clean chain enters only as the regression target; the interface-B
variant of E03 and E14 keeps the same loss but with one fixed $(\vz,t)$ per chain.

\subsection{Implicit score matching (ISM)}

Under interface~C there is no clean target, and the network minimises Hyv\"arinen's
objective on the noised views alone:
\begin{equation}
  \mathcal L_{\mathrm{ISM}}(\theta)=\E\Bigl[\tfrac12\lVert\hat\score_\theta(\vx,t)\rVert^2+\nabla_{\vx}\!\cdot\hat\score_\theta(\vx,t)\Bigr],
  \qquad
  \hat\score_\theta=-\frac{\vx-\alpha\hat m_\theta}{\Dt},
  \label{eq:ism-loss}
\end{equation}
where $\hat m_\theta$ is formed from the network output by the parameterisation in use;
the functions are \code{ism\_loss} and \code{\_score\_from\_out}.
\Cref{ch:model-selection} derives the identity
\begin{equation}
  \mathcal L_{\mathrm{ISM}}(\theta)=\tfrac12\E\lVert\hat\score_\theta-\score\rVert^2+\text{const}
\end{equation}
that makes this a score-matching loss without the score. Its cost is the divergence term,
a sum of $\nsites$ partial derivatives of the network's own output with respect to its
own input, which must be differentiated again for the gradient step
(\code{create\_graph=True}). The batches are drawn from the fixed set of one-view samples
(no redrawing, there is nothing to redraw), each with the level it was observed at.

\section{Three divergence estimators}
\label{sec:divergence}

Let $\mat{J}(\vx)=\partial\hat\score/\partial\vx$ be the $\nsites\times\nsites$ Jacobian of
the score map at $\vx$. The divergence is $\tr\mat{J}=\sum_iJ_{ii}$. Automatic
differentiation provides, for a cost comparable to one forward pass, the product
$\mat{J}^\top\vect{v}$ for any vector $\vect{v}$ (one backward pass through the network),
but not $\mat{J}$ itself. The function \code{divergence} offers three ways to obtain the
trace from such products.

\paragraph{Exact.} Take $\vect{v}=\vect{e}_i$ for $i=1,\dots,\nsites$: the $i$-th backward
pass returns row $i$ of $\mat{J}$, whose $i$-th entry is $J_{ii}$. Cost: $\nsites$
backward passes per batch, thirty-two at the paper's length, at every training step.

\paragraph{Hutchinson's estimator.} Draw one probe $\vect{v}\sim\normal0{\Id}$ and use
$\vect{v}^\top\mat{J}\vect{v}$. Since $\E[v_iv_j]=\delta_{ij}$,
$\E[\vect{v}^\top\mat{J}\vect{v}]=\sum_{ij}J_{ij}\E[v_iv_j]=\tr\mat{J}$: unbiased at the cost
of one backward pass. Its variance is $\sum_{i\ne j}J_{ij}^2+\sum_{ij}J_{ij}J_{ji}$ for a
Gaussian probe, of the order of the squared off-diagonal mass of $\mat{J}$, which for a
dense map at low noise is large; the estimator adds noise to every gradient step, and
the audit's record of E14's first screen (\Cref{ch:protocol}) shows what that noise did to
a window head trained with it.

\paragraph{Coloured probes for a banded Jacobian.} If output $i$ depends only on inputs
$j$ with $|i-j|\le r$, then $\mat{J}$ is banded with bandwidth $r$, and $p=2r+1$
deterministic probes give the trace exactly.

\begin{proposition}[Coloured trace]
\label{prop:coloured}
Let $J_{ij}=0$ whenever $|i-j|>r$, let $p=2r+1$ and let $\vect{v}^{(k)}\in\{0,1\}^{\nsites}$
have entries $v^{(k)}_j=\mathbf 1\{j\equiv k\pmod p\}$ for $k=0,\dots,p-1$. Then
$\sum_{k=0}^{p-1}\vect{v}^{(k)\top}\mat{J}\vect{v}^{(k)}=\tr\mat{J}$.
\end{proposition}
\begin{proof}
$\vect{v}^{(k)\top}\mat{J}\vect{v}^{(k)}=\sum_{i\equiv k}\sum_{j\equiv k}J_{ij}$. Two indices in the
same residue class differ by a multiple of $p=2r+1$, so either $i=j$ or $|i-j|\ge p>r$,
in which case $J_{ij}=0$. Hence the sum is $\sum_{i\equiv k}J_{ii}$, and summing over the
$p$ classes covers every $i$ once.
\end{proof}

This is the graph-colouring idea of Curtis, Powell and Reid for sparse Jacobians: sites
of the same colour do not interact, so one probe per colour suffices. The zero padding at
the ends only removes dependencies and never adds one, so the band structure holds at
the boundary sites too. For the window head the code's \code{receptive\_radius} supplies
$r$; for a dense map the function refuses the coloured mode, since a transformer's or
GRU's output at $i$ depends on every input and the proposition does not apply. Cost:
$2r+1$ backward passes, nine at radius $4$, against thirty-two for the exact loop, with
no variance. E14 trains the window head with the coloured probes, the transformer with
Hutchinson's estimator, and selects every arm's checkpoint with an exact trace
(coloured for the head, the $\nsites$-pass loop for the transformer); the test suite
checks that the coloured and exact traces agree to $10^{-5}$ and that Hutchinson's
estimator averages to them.

\begin{example}[Coloured probes at length five, radius one]
\label{ex:coloured}
With $\nsites=5$ and $r=1$, $p=3$ and the probes are
$\vect{v}^{(0)}=(1,0,0,1,0)$, $\vect{v}^{(1)}=(0,1,0,0,1)$, $\vect{v}^{(2)}=(0,0,1,0,0)$.
The Jacobian is tridiagonal. Then
$\vect{v}^{(0)\top}\mat{J}\vect{v}^{(0)}=J_{11}+J_{14}+J_{41}+J_{44}=J_{11}+J_{44}$ because
$|1-4|=3>1$; likewise $\vect{v}^{(1)\top}\mat{J}\vect{v}^{(1)}=J_{22}+J_{55}$ and
$\vect{v}^{(2)\top}\mat{J}\vect{v}^{(2)}=J_{33}$; the three sum to the trace, with three
backward passes instead of five. At $\nsites=128$ and $r=4$ the saving is nine passes
against a hundred and twenty-eight, which is what made the interface-C window head
affordable in the longer regime of E14.
\end{example}

\section{What the experiments do with the networks}
\label{sec:net-experiments}

\paragraph{E03: the screen, the confirmatory table, the matched arm and the length study.}
\Cref{ch:model-selection} describes the screen and the confirmation. The
\emph{matched} part trains the two families that carry the interface-C comparison, the
window head and the transformer, by implicit score matching on the one-view groups with
the exact divergence for the head and Hutchinson's estimator for the transformer, and
selects on the held-out implicit objective; beside each it trains the interface-B arm
(denoising loss on the clean chain with one fixed corruption). The \emph{length} part
loads each family's selected confirmatory checkpoint at $\nsites=32$ (saved under
\code{outputs/E03/ckpts/}) and scores it, unchanged, against the tree-exact score of chains
of length $64$, $128$ and $256$; the MLP is excluded as not length-agnostic. The paper's
Section~5 reports all four parts (the baseline table, the matched paragraph, the length
table) and its Appendix A3 the screen.

\paragraph{E08: one window head against the misspecified chain.}
The misspecification experiment's network arm is a radius-$2$, width-$64$ window head
trained under interface~A for $16{,}000$ steps with learning rate $10^{-3}$ and the eps
parameterisation, checkpoint on the validation bulk risk (\Cref{ch:misspecification-edgeworth}).
The paper names its configuration in the figure caption and the protocol appendix; it is
not the screen's radius-$4$ winner, and the audit asked that this be said.

\paragraph{E14: the matched comparison across regimes.}
The network arms train exactly as E03's matched winners: the window head with the
coloured-probe divergence at learning rate $10^{-3}$, the transformer with Hutchinson's
estimator at $3\cdot10^{-3}$, eps parameterisation for both, re-screened inside each
$(\corr,\nsites)$ regime over E03's candidate grid on two development seeds for $8{,}000$
steps, then confirmed for $32{,}000$; checkpoints selected on the held-out implicit
objective with the exact trace. The arms are \code{em\_bp}, \code{em\_bp\_evidence},
\code{ism\_one\_view}, \code{dsm\_one\_corruption}, \code{lmmse\_true} and
\code{isotropic}, each defined in \Cref{ch:model-selection}. The device is chosen at run
time (\code{pick\_device}), so the same code runs on the cluster's GPUs; state dictionaries
are always saved on the CPU so a checkpoint is portable.

\paragraph{E05: the networks as samplers.}
The generation experiment drives its reverse SDE (\Cref{ch:reverse-sde-generation}) with
each arm's score, and its network arms are the selected confirmatory checkpoints of E03,
loaded per cell so that the network is paired with the other arms within a seed.

\where{\begingroup\sloppy The paper, Section~5 (baselines under one recipe; the
comparison at matched information; length extrapolation), Section~1 (the list of
architectures), Appendix A3 (the recipe and the screen table) and the misspecification
figure. Code: \code{src/markov\_score/nn\_torch.py} (all architectures, \code{TrainConfig},
\code{train}, \code{dsm\_loss}, \code{ism\_loss}, \code{divergence}, \code{ism\_objective},
\code{Denoiser}), \code{tests/test\_nn\_torch.py}, \code{experiments/e03\_baselines.py},
\code{experiments/e14\_regime\_baselines.py}, \code{experiments/e08\_misspecification.py},
\code{experiments/e05\_generation.py}, \code{tools/make\_baseline\_numbers.py}.\par\endgroup}

\paragraph{References.}
\begingroup\sloppy
P.~Vincent, 2011, A connection between score matching and denoising autoencoders, \emph{Neural Computation} 23, 1661--1674.
A.~Hyv\"arinen, 2005, Estimation of non-normalized statistical models by score matching, \emph{Journal of Machine Learning Research} 6, 695--709.
M.~F.~Hutchinson, 1989, A stochastic estimator of the trace of the influence matrix for Laplacian smoothing splines, \emph{Communications in Statistics: Simulation and Computation} 18, 1059--1076.
A.~R.~Curtis, M.~J.~D.~Powell and J.~K.~Reid, 1974, On the estimation of sparse Jacobian matrices, \emph{Journal of the Institute of Mathematics and its Applications} 13, 117--119.
D.~P.~Kingma and J.~Ba, 2015, Adam: a method for stochastic optimization, \emph{International Conference on Learning Representations}.
K.~Cho, B.~van Merri\"enboer, C.~Gulcehre, D.~Bahdanau, F.~Bougares, H.~Schwenk and Y.~Bengio, 2014, Learning phrase representations using RNN encoder--decoder for statistical machine translation, \emph{EMNLP}.
A.~Vaswani et al., 2017, Attention is all you need, \emph{Advances in Neural Information Processing Systems} 30.
S.~Bai, J.~Z.~Kolter and V.~Koltun, 2018, An empirical evaluation of generic convolutional and recurrent networks for sequence modeling, arXiv:1803.01271.
O.~Ronneberger, P.~Fischer and T.~Brox, 2015, U-Net: convolutional networks for biomedical image segmentation, \emph{MICCAI}.
O.~Press, N.~A.~Smith and M.~Lewis, 2022, Train short, test long: attention with linear biases enables input length extrapolation, \emph{International Conference on Learning Representations}.\par\endgroup
